\chapterhead{54}{ORR-15}{Supervisory Guidance on Model Risk Management}{1}{3}

\lohead{54}{a}{Describe model risk and explain how it can arise in the
implementation of a model.}

\begin{ordefbox}{Model risk}
Model risk pertains to the potential for adverse outcomes arising from using
flawed models in financial decision-making, leading to financial or reputational
losses.
\end{ordefbox}

\orsub{Components of a model}

\begin{center}
\begin{tikzpicture}[font=\fontsize{8.2}{10.4}\selectfont,
  st/.style={text width=4.2cm,align=center,rounded corners=3pt,inner sep=6pt,
             text=white,font=\sffamily\bfseries\fontsize{8.8}{10.6}\selectfont},
  sb/.style={text width=4.2cm,align=center,rounded corners=3pt,inner sep=7pt,
             fill=paleblue,draw=palenavy,line width=0.6pt,text=navy}]
\foreach \i/\x/\c/\t/\d in {%
  1/-5.6/navy/{(1) Information inputs}/{Data and assumptions},
  2/0/teal/{(2) Processing}/{Conversion of inputs to estimates},
  3/5.6/forest/{(3) Reporting}/{Conversion of estimates to applied information}}
{
  \node[st,fill=\c,anchor=north] (t\i) at (\x,0) {\t};
  \node[sb,anchor=north] (b\i) at (t\i.south) {\d};
}
\foreach \a/\b in {1/2,2/3}
  \draw[-{Stealth[length=5pt]},navy!35,line width=1.4pt]
    ($(t\a.east)+(0.05,0)$) -- ($(t\b.west)+(-0.05,0)$);
\end{tikzpicture}
\end{center}
\orfigcap{The three components of a model.}

\orsub{Sources of model risk}

\begin{itemize}
\item \textbf{Model errors:} faulty outputs due to errors in design,
implementation, or simplifying assumptions.
\item \textbf{Misuse:} using models improperly, outside their intended scope, or
in changed market conditions.
\end{itemize}

% ---------------------------------------------------------------------
\lohead{54}{b}{Describe elements of an effective model risk management process.}

\orsub{Effective model risk management elements}

\begin{lettered}
\item Analyse the complexity, uncertainty, and potential impact of models to
determine the level of risk management needed.
\item Have a separate team with strong technical skills objectively evaluate the
model and its assumptions; this should be independent of model development.
\item Use model outputs along with other information to make informed decisions.
\item Continuously track model performance and make adjustments as needed,
supported by management.
\item Include restrictions on model use, performance analysis, and continual
model calibration.
\end{lettered}

% ---------------------------------------------------------------------
\lohead{54}{c}{Explain best practices for the development and implementation of
models.}

\orsub{Best practices for model development and implementation}

\begin{lettered}
\item Begin with a clear statement of the model's purpose, aligned with its
eventual use. Thoroughly describe the supporting information, noting strengths
and weaknesses.
\item Ensure the model is based on correct mathematical theories and compare it
with alternative models for a ``reasonability check''.
\item Ensure data and assumptions are relevant and robust. Document any data
proxies or adjustments clearly.
\item Test the model for precision, strength, and consistency across normal and
stressed market conditions. Document the testing process comprehensively.
\item Use multiple tests to avoid Type I and Type II errors, considering both
quantitative and qualitative aspects.
\item Implement robust internal controls involving both individuals and
technology to support effective model development and implementation.
\end{lettered}

% ---------------------------------------------------------------------
\lohead{54}{d}{Describe elements of a strong model validation process and
challenges to an effective validation process.}

\orsub{Model validation process}

\begin{orkeybox}
Model validation is a series of checks to ensure models perform as intended. It
is crucial to have a separate team, ideally independent from the model
developers, to conduct validation.
\end{orkeybox}

\orsubb{1) Evaluation of conceptual soundness}

This is a quality check on the entire model. It involves reviewing documentation,
test results, and the model's development process. Key aspects include:

\begin{lettered}
\item Reviewing major assumptions and how they impact outputs.
\item Assessing the data's relevance and representativeness for the model's
purpose.
\item Performing sensitivity analysis to see how input changes affect outputs
(single and multiple variables).
\item Stress testing the model with extreme values to see if it breaks.
\item Validating non-quantitative aspects (subjective judgments), ensuring proper
backup and documentation.
\end{lettered}

\orsubb{2) Ongoing monitoring}

This checks if the model needs adjustments or replacement after implementation.
It involves:

\begin{lettered}
\item \textbf{Process verification:} ensuring data inputs are error-free and
complete, and computer code is accurate.
\item \textbf{Benchmarking:} comparing the model's outputs with similar data or
models to identify significant variances.
\item Monitoring whether the model remains effective or needs updates.
\end{lettered}

\orsubb{3) Outcomes analysis}

This examines the model's outputs compared to actual results. It includes:

\begin{lettered}
\item \textbf{Precision analysis:} evaluating the accuracy of the model's
estimates, through quantitative and qualitative methods.
\item \textbf{Parallel outcomes analysis:} comparing estimates from the original
and amended models with actual outcomes.
\item \textbf{Backtesting:} analyse variances between actual results and model
estimates over different periods and economic conditions. Backtesting can be
complex due to data limitations and the need for additional testing over shorter
periods for long forecasts.
\item If significant problems are found, the model may need amendments,
reconstruction, or even replacement. Revalidation is required after major
changes.
\item \textbf{Early warning signals:} using real-time metrics to monitor model
performance.
\end{lettered}

\orsubb{4) Vendor and other third-party products}

\begin{lettered}
\item Limited access to internal workings of the model may necessitate more focus
on sensitivity analysis and benchmarking during validation.
\item Require vendors to provide detailed information on their models for
effective validation.
\end{lettered}
